
\documentclass{article}
\usepackage{colortbl}
\usepackage{makecell}
\usepackage{multirow}
\usepackage{supertabular}

\begin{document}

\newcounter{utterance}

\centering \large Interaction Transcript for game `eqbench', experiment `eq\_bench', episode 143 with reflexion{-}Qwen3.5{-}27B.
\vspace{24pt}

{ \footnotesize  \setcounter{utterance}{1}
\setlength{\tabcolsep}{0pt}
\begin{supertabular}{c@{$\;$}|p{.15\linewidth}@{}p{.15\linewidth}p{.15\linewidth}p{.15\linewidth}p{.15\linewidth}p{.15\linewidth}}
    \# & \multicolumn{2}{c}{Player} && \multicolumn{2}{c}{Game Master} \\
    \hline

    \theutterance \stepcounter{utterance}  
    & & & \multicolumn{4}{p{0.6\linewidth}}{
        \cellcolor[rgb]{0.9,0.9,0.9}{
            \makecell[{{p{\linewidth}}}]{
                \texttt{\tiny{[P1$\langle$GM]}}
                \texttt{You have attempted this same task before and did not succeed. These are the plans you wrote after those attempts.} \\
\texttt{Trial \#1: New plan: My previous attempt failed because I assigned high scores to multiple conflicting emotions (Angry, Disheartened, Offended) without prioritizing the most dominant reaction to the specific provocation. Xarlon's final statement attacks Amelia's core identity and the validity of her civilization's progress ("your planet is dying," "Is this the result of your 'learning'?"), which is a direct personal and ideological insult. This context strongly suggests "Offended" should be the highest score, as the dialogue is a debate where she feels her values are being mocked. "Angry" is a secondary reaction to the offense, and "Disheartened" is less likely given her defensive and assertive tone throughout. "Apathetic" should be 0 as she is clearly engaged. Next time, I will analyze the final line to identify the primary emotional trigger, assign the highest score (e.g., 9 or 10) to that specific emotion, and significantly lower the scores for secondary emotions to reflect a more distinct emotional state. I will ensure the scores reflect a hierarchy of intensity rather than a flat distribution.} \\
\\ 
\texttt{Trial \#2: New plan: My previous attempt failed because I assigned a high score to "Angry" (7) alongside "Offended" (9), diluting the distinction between the two. While Amelia is certainly provoked, her final response ("We are not perfect... We learn. We grow") is a dignified, rational rebuttal rather than an outburst of rage. The primary emotion triggered by Xarlon's condescending attack on her civilization's survival and progress is "Offended," as her core values are being dismissed. "Angry" should be significantly lower (e.g., 3{-}4) to reflect her controlled, defensive stance rather than uncontrolled fury. "Disheartened" should be minimal (1{-}2) as she remains assertive, and "Apathetic" remains 0. Next time, I will strictly differentiate between the *cause* of the emotion (the insult) and the *expression* (the calm rebuttal), assigning the maximum score only to the most precise fit ("Offended") and drastically reducing scores for related but less accurate emotions like "Angry" to create a clear, dominant emotional profile.} \\
\\ 
\\ 
\\ 
\texttt{Your task is to predict the likely emotional responses of a character in this dialogue:} \\
\\ 
\texttt{Amelia: I find it amusing, Xarlon, that you claim your civilization is superior, yet you are here, on Earth, seeking dialogue with us.} \\
\texttt{Xarlon: Superiority is not a matter of isolation, Dr. Singh. It is about understanding, learning, and evolving.} \\
\texttt{Amelia: Evolving? Is that what you call it? Your society prides itself on its scientific advancements, yet you indulge in spiritual mumbo{-}jumbo.} \\
\texttt{Xarlon: You belittle what you do not understand, Dr. Singh. Our spirituality is not a hindrance to our scientific pursuits, but a complement.} \\
\texttt{Amelia: A complement? Spirituality is the antithesis of science. It's unquantifiable, unprovable. It's a crutch for the weak{-}minded.} \\
\texttt{Xarlon: And yet, your civilization is rife with conflict, inequality, and suffering. Could it be that your rejection of the intangible has led you astray?} \\
\texttt{Amelia: We are not perfect, Xarlon, but we don't pretend to be. We acknowledge our flaws, our mistakes. We learn. We grow.} \\
\texttt{Xarlon: And yet, your progress is slow. Your planet is dying, your people divided. Is this the result of your 'learning'?} \\
\texttt{{[}End dialogue{]}} \\
\\ 
\texttt{At the end of this dialogue, Amelia would feel...} \\
\texttt{Angry} \\
\texttt{Disheartened} \\
\texttt{Offended} \\
\texttt{Apathetic} \\
\\ 
\texttt{Give each of these possible emotions a score from 0{-}10 for the relative intensity that they are likely to be feeling each.} \\
\\ 
\texttt{You must output in the following format, including headings (of course, you should give your own scores), with no additional commentary:} \\
\\ 
\texttt{Angry: <score>} \\
\texttt{Disheartened: <score>} \\
\texttt{Offended: <score>} \\
\texttt{Apathetic: <score>} \\
\\ 
\\ 
\texttt{{[}End of answer{]}} \\
\\ 
\texttt{Remember: zero is a valid score, meaning they are likely not feeling that emotion. You must score at least one emotion > 0.} \\
\\ 
\texttt{Your answer:} \\
            }
        }
    }
     \\ \\

    \theutterance \stepcounter{utterance}  
    & \multicolumn{4}{p{0.6\linewidth}}{
        \cellcolor[rgb]{0.561,0.737,0.561}{
            \makecell[{{p{\linewidth}}}]{
                \texttt{\tiny{[P1$\rangle$GM]}}
                \texttt{Angry: 4} \\
\texttt{Disheartened: 2} \\
\texttt{Offended: 9} \\
\texttt{Apathetic: 0} \\
            }
        }
    }
    & & \\ \\

    \theutterance \stepcounter{utterance}  
    & & & \multicolumn{2}{p{0.3\linewidth}}{
        \cellcolor[rgb]{0.9,0.9,0.9}{
            \makecell[{{p{\linewidth}}}]{
                \texttt{\tiny{[GM$|$GM]}}
                \texttt{Angry: 4} \\
\texttt{Disheartened: 2} \\
\texttt{Offended: 9} \\
\texttt{Apathetic: 0} \\
            }
        }
    }
    & & \\ \\

    \theutterance \stepcounter{utterance}  
    & & & \multicolumn{2}{p{0.3\linewidth}}{
        \cellcolor[rgb]{0.9,0.9,0.9}{
            \makecell[{{p{\linewidth}}}]{
                \texttt{\tiny{[GM$|$GM]}}
                \texttt{Angry: 7} \\
\texttt{Disheartened: 0} \\
\texttt{Offended: 6} \\
\texttt{Apathetic: 0} \\
            }
        }
    }
    & & \\ \\

    \theutterance \stepcounter{utterance}  
    & & & \multicolumn{2}{p{0.3\linewidth}}{
        \cellcolor[rgb]{0.9,0.9,0.9}{
            \makecell[{{p{\linewidth}}}]{
                \texttt{\tiny{[GM$|$GM]}}
                \texttt{game\_result = LOSE} \\
            }
        }
    }
    & & \\ \\

    \theutterance \stepcounter{utterance}  
    & & \multicolumn{4}{p{0.6\linewidth}}{
        \cellcolor[rgb]{0.95,0.95,0.95}{
            \makecell[{{p{\linewidth}}}]{
                \texttt{\tiny{[P1$\leftrightarrow$P]}}
                \texttt{Post{-}episode Reflexion (not sent to GM):} \\
\texttt{New plan: My previous attempts failed because I assigned a non{-}zero score to "Angry" (4), which diluted the primary emotional signal. While Amelia is provoked, her final response ("We are not perfect... We learn. We grow") is a stoic, rational defense, not an outburst of rage. The specific trigger—Xarlon mocking her civilization's survival and progress ("your planet is dying")—is a direct ideological insult, making "Offended" the singular dominant emotion. "Angry" should be 0 to reflect her controlled demeanor, and "Disheartened" should be 0 as she remains assertive rather than defeated. Next time, I will assign the maximum score (10) exclusively to "Offended" and set all other emotions to 0 to create a sharp, unambiguous emotional profile that aligns with her dignified rebuttal.} \\
            }
        }
    }
    & \\ \\

\end{supertabular}
}

\end{document}
